%% -----------------------------------------------
%% Memoria del Trabajo de Fin de Grado
%% Grado en Ingeniería Informática — EPS UIB
%% Codificación: UTF-8
%% -----------------------------------------------

\documentclass[spanish,GINF]{TFGEPSUIB}

% Expresiones matemáticas
\usepackage{amsmath,mathtools}

% Fragmentos de código fuente con resaltado de sintaxis
\usepackage{listings}

% Cajas destacadas para notas, advertencias, etc.
\usepackage{tcolorbox}

% Lista de acrónimos
\usepackage[printonlyused]{acronym}

% Unidades de medida y números en tablas
\usepackage{siunitx}

% Definición del lenguaje TypeScript para listings
\lstdefinelanguage{TypeScript}{
  keywords={const,let,var,function,return,if,else,for,while,do,switch,
            case,break,continue,new,typeof,instanceof,class,extends,
            export,import,default,async,await,try,catch,finally,throw,
            null,undefined,true,false,interface,type,enum,namespace,
            readonly,public,private,protected,from,of,in,as},
  keywordstyle=\color{blue}\bfseries,
  ndkeywords={string,number,boolean,void,any,never,object,unknown,Promise},
  ndkeywordstyle=\color{teal},
  sensitive=true,
  comment=[l]{//},
  morecomment=[s]{/*}{*/},
  commentstyle=\color{gray}\itshape,
  stringstyle=\color{orange!80!black},
  morestring=[b]",
  morestring=[b]',
  morestring=[b]`,
}

\lstdefinelanguage{SQL}{
  keywords={SELECT,FROM,WHERE,JOIN,LEFT,RIGHT,INNER,ON,INSERT,INTO,VALUES,
            UPDATE,SET,DELETE,CREATE,TABLE,ALTER,ADD,COLUMN,DROP,INDEX,
            NOT,NULL,DEFAULT,PRIMARY,KEY,FOREIGN,REFERENCES,UNIQUE,
            IF,EXISTS,AND,OR,IN,IS,LIKE,BETWEEN,ORDER,BY,GROUP,HAVING,
            LIMIT,OFFSET,WITH,AS,RETURNS,LANGUAGE,SECURITY,DEFINER,
            POLICY,ENABLE,ROW,LEVEL,FOR,ALL,USING,CHECK,TRIGGER,FUNCTION},
  keywordstyle=\color{blue}\bfseries,
  sensitive=false,
  comment=[l]{--},
  morecomment=[s]{/*}{*/},
  commentstyle=\color{gray}\itshape,
  stringstyle=\color{orange!80!black},
  morestring=[b]',
}

% Configuración global de listings
\lstset{
  basicstyle=\ttfamily\small,
  numbers=left,
  numberstyle=\tiny\color{gray},
  stepnumber=1,
  numbersep=8pt,
  breaklines=true,
  breakatwhitespace=false,
  frame=single,
  framerule=0.4pt,
  rulecolor=\color{gray!40},
  backgroundcolor=\color{gray!5},
  captionpos=b,
  tabsize=2,
  showspaces=false,
  showstringspaces=false,
  literate={á}{{\'a}}1 {é}{{\'e}}1 {í}{{\'i}}1 {ó}{{\'o}}1 {ú}{{\'u}}1
           {Á}{{\'A}}1 {É}{{\'E}}1 {Í}{{\'I}}1 {Ó}{{\'O}}1 {Ú}{{\'U}}1
           {ñ}{{\~n}}1 {Ñ}{{\~N}}1,
}

% hyperref debe cargarse al final del preámbulo
\usepackage[backref, colorlinks=true, allcolors=black]{hyperref}

%% -----------------------------------------------
%% Datos de la portada
%% -----------------------------------------------

% Título del TFG
\title{Diseño e implementación de una plataforma web para la venta de productos digitales con gestión de licencias y pagos integrados}

% Nombre completo del autor
\author{[Nombre completo del autor]}

% Nombre del tutor/tutora
\tutor{[Nombre del tutor/tutora]}

% Año académico
\date{Año académico 2025-26}

% Palabras clave (separadas por comas)
\paraulesclau{plataforma web, productos digitales, gestión de licencias, comercio electrónico, Next.js, Supabase, Stripe, TypeScript}

% Descomentar si NO se autoriza la publicación en abierto:
%\autorfalse
%\tutorfalse

%% Para compilar solo algunos capítulos mientras se redacta:
%\includeonly{Introduccion}

\begin{document}

\portada
\portadainterior
\frontmatter

%% Agradecimientos (opcional — descomentar para activar)
%\cleartorecto \thispagestyle{empty}
%\begin{agraiments}
%Texto de agradecimientos.
%\end{agraiments}

% Índice de contenidos
\cleartorecto \tableofcontents

% Lista de figuras (opcional)
%\cleartorecto \listoffigures

% Lista de tablas (opcional)
%\cleartorecto \listoftables

% Lista de acrónimos
\include{Acronimos}

% Resumen en español + Abstract en inglés
%!TeX root=MemoriaTFG.tex

\chapter{Resumen}

% 250-500 palabras. Debe ser completamente autónomo: el lector entiende el
% trabajo leyendo solo esta sección. Sin acrónimos sin definir, sin citas.
%
% Estructura:
%   1. ¿Qué necesidad cubre la plataforma? ¿Qué carencias del mercado justifican su desarrollo?
%   2. ¿Cómo se ha abordado el desarrollo? (metodología, arquitectura elegida)
%   3. ¿Qué se ha construido? ¿Qué funcionalidades ofrece? ¿Qué métricas se han alcanzado?
%   4. ¿Se han cumplido los objetivos? ¿Qué valor aporta el trabajo?

[Redactar el resumen en español aquí.]


\section*{Abstract}

% Same content as above, translated to English.
% Both versions are required for submission to the UIB institutional repository.

[Write the abstract in English here.]


%% -----------------------------------------------
%% Cuerpo del trabajo
%% -----------------------------------------------
\mainmatter\pagestyle{ruled}

\include{Introduccion}
\include{EstadoArte}
\include{Diseno}
\include{Implementacion}
\include{Pruebas}
\include{Resultados}
\include{Conclusiones}

%% -----------------------------------------------
%% Anexos
%% -----------------------------------------------
\appendix
\include{Anexos}

%% -----------------------------------------------
%% Bibliografía
%% -----------------------------------------------
\backmatter
\bibliographystyle{IEEEtran}
\bibliography{Bibliografia}

\end{document}
